\documentclass[conference]{IEEEtran}
\IEEEoverridecommandlockouts

\usepackage{cite}
\usepackage{amsmath,amssymb,amsfonts}
\usepackage{graphicx}
\usepackage{booktabs}
\usepackage[utf8]{inputenc}
\usepackage[T1]{fontenc}
\usepackage[french]{babel}
\usepackage{subcaption}
\usepackage{tabularx}
\usepackage{url}
\usepackage{float}
\usepackage[hidelinks]{hyperref}
\usepackage{placeins}
\usepackage{listings}
\usepackage{tikz}
\usepackage{pgfplots}
\pgfplotsset{compat=1.18}
\usetikzlibrary{shapes,arrows.meta,positioning,fit,calc}

\lstdefinestyle{p4log}{
  basicstyle=\ttfamily\scriptsize,
  breaklines=true,
  columns=fullflexible,
  frame=single,
  framesep=3pt,
  xleftmargin=1pt,
  xrightmargin=1pt,
  showstringspaces=false,
  literate={\#}{{\#}}1
}

\newcommand{\toolname}{\textsc{Argos}}
\newcommand{\beepbeep}{\textsc{BeepBeep}}

% =============================================================================
% Figures TikZ / pgfplots pour Argos (remplace les PNG flous)
% Prérequis dans le préambule :
%   \usepackage{tikz}
%   \usepackage{pgfplots}
%   \pgfplotsset{compat=1.18}
%   \usetikzlibrary{shapes,arrows.meta,positioning,fit,calc}
% Usage : % =============================================================================
% Figures TikZ / pgfplots pour Argos (remplace les PNG flous)
% Prérequis dans le préambule :
%   \usepackage{tikz}
%   \usepackage{pgfplots}
%   \pgfplotsset{compat=1.18}
%   \usetikzlibrary{shapes,arrows.meta,positioning,fit,calc}
% Usage : % =============================================================================
% Figures TikZ / pgfplots pour Argos (remplace les PNG flous)
% Prérequis dans le préambule :
%   \usepackage{tikz}
%   \usepackage{pgfplots}
%   \pgfplotsset{compat=1.18}
%   \usetikzlibrary{shapes,arrows.meta,positioning,fit,calc}
% Usage : \input{figures_tikz} puis \figXX... dans le document
% Données extraites du pipeline sur log.txt (run juin 2026) :
% counters, gauges, histogrammes, summaries, alertes JSON.
\newcommand{\RushFenetres}{24} % len(rush_collectif) dans perforce_events_alertes.json

% Coordonnées sûres pour pgfplots (éviter accents / math dans symbolic coords)
\pgfplotsset{
  /pgf/number format/.cd,
  set thousands separator={\,},
}

% --- fig00 : architecture pipeline -------------------------------------------
\newcommand{\figZeroArchitecture}{%
\begin{figure}[!t]
  \centering
  \resizebox{\linewidth}{!}{%
  \begin{tikzpicture}[
    font=\scriptsize,
    box/.style={draw, rounded corners=2pt, align=center, inner sep=4pt,
                minimum width=3.4cm, fill=white},
    small/.style={box, minimum width=3.0cm, font=\tiny},
    arr/.style={-{Stealth[length=2mm]}, thick, blue!60!black},
    darr/.style={-{Stealth[length=2mm]}, thick, gray!55},
    lbl/.style={font=\tiny, text=gray}
  ]
    \node[box, fill=blue!8, minimum width=4.2cm] (log) {\textbf{log.txt}\\{\tiny entrée (${\sim}$1,7~Go)}};
    \node[box, below=0.45cm of log] (parse) {\textbf{LecteurLog}\\{\tiny bloc = cmd + completed + \texttt{---} (pid)}};
    \node[box, below=0.45cm of parse] (dedup) {\textbf{Déduplication}\\{\tiny timestamp|pid|cmd|args|fichiers}};
    \node[small, below left=0.75cm and 0.15cm of dedup, fill=orange!10] (ds)
      {\textbf{Export v3}\\{\tiny JSON / CSV (27 col.)}};
    \node[box, below right=0.55cm and -0.1cm of dedup, fill=blue!6,
          minimum width=3.6cm] (bb)
      {\textbf{BeepBeep 3.13}\\{\tiny \texttt{QueueSource} $\rightarrow$ processeurs}};
    \node[small, below left=0.55cm and 0.05cm of bb, fill=green!8] (p1)
      {\textbf{Phase 1}\\{\tiny Counter, Gauge,\\Histogram, Summary}};
    \node[small, below right=0.55cm and 0.05cm of bb, fill=red!6] (p2)
      {\textbf{Phase 2 (CEP)}\\{\tiny \texttt{DetecteurPatterns}\\edit$\rightarrow$revert · apathie · rush}};
    \draw[arr] (log) -- (parse);
    \draw[arr] (parse) -- (dedup);
    \draw[darr] (dedup.south) -- ++(0,-0.22) -| (ds.north);
    \draw[arr] (dedup.south) -- ++(0,-0.22) -| (bb.north);
    \draw[arr] (bb.south) -- ++(0,-0.18) -| (p1.north);
    \draw[arr] (bb.south) -- ++(0,-0.18) -| (p2.north);
    \node[lbl, below=0.12cm of ds] {\texttt{perforce\_events\_v3}};
    \node[lbl, below=0.12cm of p1] {4 JSON OpenMetrics};
    \node[lbl, below=0.12cm of p2] {\texttt{perforce\_events\_alertes.json}};
  \end{tikzpicture}%
  }
  \caption{Pipeline en flux~: parse $\rightarrow$ dédup $\rightarrow$ export v3 et analyse \beepbeep{} (phases~1 et~2).}
  \label{fig:arch}
\end{figure}}

% --- fig01 : compteurs -------------------------------------------------------
\newcommand{\figUnCompteurs}{%
\begin{figure}[!t]
  \centering
  \begin{tikzpicture}
    \begin{axis}[
      width=\linewidth,
      height=4.2cm,
      ybar,
      bar width=14pt,
      ymin=0,
      symbolic x coords={Evts,Edits,Submits},
      xticklabels={Événements,Edits,Submits},
      xtick=data,
      ylabel={Nombre},
      tick label style={font=\scriptsize},
      label style={font=\scriptsize},
      nodes near coords,
      nodes near coords style={font=\tiny},
      enlarge x limits=0.2,
      grid=major,
      grid style={gray!25},
    ]
      \addplot[fill=blue!55!black] coordinates {
        (Evts,16262) (Edits,4069) (Submits,1224)
      };
    \end{axis}
  \end{tikzpicture}
  \caption{Compteurs globaux (phase~1).}
  \label{fig:compteurs}
\end{figure}}

% --- fig02 : submits par heure (points clés 27/04) ---------------------------
\newcommand{\figDeuxSubmitsHeure}{%
\begin{figure}[!t]
  \centering
  \begin{tikzpicture}
    \begin{axis}[
      width=\linewidth,
      height=4cm,
      xlabel={Heure (27/04)},
      ylabel={Submits},
      tick label style={font=\scriptsize},
      label style={font=\scriptsize},
      grid=major,
      grid style={gray!25},
      xmin=9, xmax=21, ymin=0, ymax=120,
      xtick={9,10,11,12,13,14,15,16,17,18,19,20,21},
      xticklabels={9h,10h,11h,12h,13h,14h,\textbf{15h},16h,17h,18h,19h,20h,21h},
    ]
      \addplot[thick, blue!70!black, mark=*] coordinates {
        (9,25) (10,48) (11,86) (12,63) (13,92) (14,95) (15,107)
        (16,81) (17,96) (18,75) (19,75) (20,81) (21,60)
      };
      \node[font=\tiny, anchor=south west] at (axis cs:15,107) {pic 107};
    \end{axis}
  \end{tikzpicture}
  \caption{Soumissions par heure (pic 107 à 15~h).}
  \label{fig:submits_heure}
\end{figure}}

% --- fig03 : jauge fichiers / utilisateurs ouverts ---------------------------
\newcommand{\figTroisGaugeOuverts}{%
\begin{figure}[!t]
  \centering
  \begin{tikzpicture}
    \begin{axis}[
      width=\linewidth,
      height=4cm,
      xlabel={Heure (27/04)},
      ylabel={Pic horaire},
      tick label style={font=\scriptsize},
      label style={font=\scriptsize},
      legend style={font=\tiny, at={(0.02,0.98)}, anchor=north west},
      grid=major,
      grid style={gray!25},
      xmin=9, xmax=20, ymin=0, ymax=520,
      xtick={10,11,12,13,14,15,16,17,18,19},
      xticklabels={10h,11h,12h,13h,14h,\textbf{15h},16h,17h,18h,19h},
    ]
      \addplot[thick, green!50!black, mark=square*] coordinates {
        (10,245) (11,311) (12,347) (13,334) (14,343) (15,294)
        (16,500) (17,286) (18,263) (19,317)
      };
      \addlegendentry{Fichiers ouverts}
      \addplot[thick, orange!80!black, mark=triangle*] coordinates {
        (10,53) (11,54) (12,64) (13,74) (14,79) (15,78)
        (16,75) (17,71) (18,72) (19,74)
      };
      \addlegendentry{Utilisateurs}
      \node[font=\tiny] at (axis cs:16,500) [above] {500};
      \node[font=\tiny] at (axis cs:14,79) [above] {79};
    \end{axis}
  \end{tikzpicture}
  \caption{Jauge — fichiers et utilisateurs en edit (pics horaires).}
  \label{fig:gauge_ouverts}
\end{figure}}

% --- fig02+03 combinées (fig:activite) ---------------------------------------
\newcommand{\figActivite}{%
\begin{figure}[!t]
  \centering
  \begin{subfigure}{\linewidth}
    \centering
    \begin{tikzpicture}
      \begin{axis}[
        width=\linewidth, height=3.2cm,
        xlabel={Heure (27/04)}, ylabel={Submits},
        tick label style={font=\tiny}, label style={font=\scriptsize},
        grid=major, grid style={gray!25},
        xmin=9, xmax=21, ymin=0, ymax=120,
        xtick={9,10,11,12,13,14,15,16,17,18,19,20,21},
        xticklabels={9h,10h,11h,12h,13h,14h,\textbf{15h},16h,17h,18h,19h,20h,21h},
      ]
        \addplot[thick, blue!70!black, mark=*] coordinates {
          (9,25) (10,48) (11,86) (12,63) (13,92) (14,95) (15,107)
          (16,81) (17,96) (18,75) (19,75) (20,81) (21,60)
        };
      \end{axis}
    \end{tikzpicture}
    \caption{Soumissions/heure (pic 107, 15~h).}
    \label{fig:submits_heure}
  \end{subfigure}\\[0.4em]
  \begin{subfigure}{\linewidth}
    \centering
    \begin{tikzpicture}
      \begin{axis}[
        width=\linewidth, height=3.2cm,
        xlabel={Heure (27/04)}, ylabel={Fichiers},
        tick label style={font=\tiny}, label style={font=\scriptsize},
        grid=major, grid style={gray!25},
        xmin=9, xmax=20, ymin=0, ymax=520,
        xtick={10,11,12,13,14,15,16,17,18,19},
        xticklabels={10h,11h,12h,13h,14h,15h,\textbf{16h},17h,18h,19h},
      ]
        \addplot[thick, green!50!black, mark=square*] coordinates {
          (10,245) (11,311) (12,347) (13,334) (14,343) (15,294)
          (16,500) (17,286) (18,263) (19,317)
        };
      \end{axis}
    \end{tikzpicture}
    \caption{Fichiers ouverts (pic 500).}
    \label{fig:gauge_ouverts}
  \end{subfigure}
  \caption{Activité sur 24~h autour de l'échéance.}
  \label{fig:activite}
\end{figure}}

% --- fig04 : histogramme durées (16 232 obs., deltas JSON) --------------------
\newcommand{\figQuatreHistogramme}{%
\begin{figure}[!t]
  \centering
  \begin{tikzpicture}
    \begin{axis}[
      width=\linewidth,
      height=4.5cm,
      ybar,
      bar width=5pt,
      ymin=0,
      ymax=6500,
      ylabel={Obs.},
      xlabel={Bucket (s)},
      tick label style={font=\tiny, rotate=45, anchor=east},
      label style={font=\scriptsize},
      symbolic x coords={%
        le005,le01,le025,le05,le1,le2,le5,le10,le30,le60,le120,le300,le600,le1200},
      xticklabels={.05,.1,.25,.5,1,2,5,10,30,60,120,300,600,1200},
      xtick=data,
      enlarge x limits=0.02,
      grid=major,
      grid style={gray!25},
      nodes near coords={},
    ]
      % deltas exacts (cumulatif → par bucket), total 16 232
      \addplot[fill=blue!50] coordinates {
        (le005,2180) (le01,6114) (le025,6202) (le05,584) (le1,301) (le2,160) (le5,189)
        (le10,87) (le30,136) (le60,81) (le120,65) (le300,48) (le600,38) (le1200,27)
      };
    \end{axis}
  \end{tikzpicture}
  \caption{Histogramme des durées (16\,232 obs., 15 buckets).}
  \label{fig:durees}
\end{figure}}

% --- fig05 : quantiles par commande ------------------------------------------
\newcommand{\figCinqQuantiles}{%
\begin{figure}[!t]
  \centering
  \begin{tikzpicture}
    \begin{axis}[
      width=\linewidth,
      height=4.2cm,
      ybar,
      bar width=4pt,
      ymin=0,
      ymax=42,
      ylabel={Secondes},
      symbolic x coords={edit,submit,add,resolve},
      xtick=data,
      legend style={font=\tiny, at={(0.98,0.98)}, anchor=north east},
      tick label style={font=\scriptsize},
      label style={font=\scriptsize},
      grid=major,
      grid style={gray!25},
      enlarge x limits=0.15,
    ]
      \addplot[fill=blue!40] coordinates {(edit,0.110) (submit,0.125) (add,0.125) (resolve,0.078)};
      \addlegendentry{p50}
      \addplot[fill=orange!60] coordinates {(edit,0.204) (submit,0.250) (add,0.452) (resolve,0.094)};
      \addlegendentry{p95}
      \addplot[fill=red!50] coordinates {(edit,0.698) (submit,0.821) (add,39.349) (resolve,0.172)};
      \addlegendentry{p99}
    \end{axis}
  \end{tikzpicture}
  \caption{Quantiles Summary par commande (p50, p95, p99).}
  \label{fig:quantiles}
\end{figure}}

% --- fig06 : zoom sync (échelle log) -----------------------------------------
\newcommand{\figSixQuantilesSync}{%
\begin{figure}[!t]
  \centering
  \begin{tikzpicture}
    \begin{axis}[
      width=0.85\linewidth,
      height=4cm,
      ybar,
      bar width=18pt,
      ymin=0.1,
      ymax=2000,
      ymode=log,
      log basis y=10,
      ylabel={Secondes (log)},
      symbolic x coords={p50,p95,p99},
      xtick=data,
      tick label style={font=\scriptsize},
      label style={font=\scriptsize},
      grid=major,
      grid style={gray!25},
      nodes near coords={\pgfmathprintnumber\pgfplotspointmeta},
      nodes near coords style={font=\tiny, /pgf/number format/fixed},
      enlarge x limits=0.3,
    ]
      \addplot[fill=purple!45] coordinates {(p50,0.188) (p95,104.65) (p99,1101.74)};
    \end{axis}
  \end{tikzpicture}
  \caption{Quantiles \texttt{user-sync} (échelle logarithmique).}
  \label{fig:quantiles_sync}
\end{figure}}

% --- fig05+06 combinées ------------------------------------------------------
\newcommand{\figQuantilesCombine}{%
\begin{figure}[!t]
  \centering
  \begin{subfigure}{0.48\linewidth}
    \centering
    \resizebox{\linewidth}{!}{%
    \begin{tikzpicture}
      \begin{axis}[
        width=6cm, height=4cm, ybar, bar width=4pt,
        ymin=0, ymax=42, ylabel={s},
        symbolic x coords={edit,submit,add,res},
        xtick=data,
        tick label style={font=\tiny},
        legend style={font=\tiny, legend columns=3, at={(0.5,1.12)}, anchor=south},
        enlarge x limits=0.12,
      ]
        \addplot[fill=blue!40] coordinates {(edit,0.110) (submit,0.125) (add,0.125) (res,0.078)};
        \addlegendentry{p50}
        \addplot[fill=orange!60] coordinates {(edit,0.204) (submit,0.250) (add,0.452) (res,0.094)};
        \addlegendentry{p95}
        \addplot[fill=red!50] coordinates {(edit,0.698) (submit,0.821) (add,39.349) (res,0.172)};
        \addlegendentry{p99}
      \end{axis}
    \end{tikzpicture}}
    \caption{Par commande.}
    \label{fig:quantiles}
  \end{subfigure}\hfill
  \begin{subfigure}{0.48\linewidth}
    \centering
    \resizebox{\linewidth}{!}{%
    \begin{tikzpicture}
      \begin{axis}[
        width=5cm, height=4cm, ybar, bar width=12pt,
        ymin=0.1, ymax=2000, ymode=log, log basis y=10,
        ylabel={s (log)},
        symbolic x coords={p50,p95,p99},
        xtick=data,
        tick label style={font=\tiny},
        nodes near coords={\pgfmathprintnumber\pgfplotspointmeta},
        nodes near coords style={font=\tiny},
        enlarge x limits=0.25,
      ]
        \addplot[fill=purple!45] coordinates {(p50,0.188) (p95,104.65) (p99,1101.74)};
      \end{axis}
    \end{tikzpicture}}
    \caption{Zoom \texttt{sync}.}
    \label{fig:quantiles_sync}
  \end{subfigure}
  \caption{Summary phase~1 (p50, p95, p99).}
\end{figure}}

% --- fig07 : alertes phase 2 (diagramme en barres, plus lisible que tarte) ---
\newcommand{\figSeptAlertes}{%
\begin{figure}[!t]
  \centering
  \begin{tikzpicture}
    \begin{axis}[
      width=\linewidth,
      height=3.8cm,
      xbar,
      bar width=10pt,
      xmin=0,
      symbolic y coords={Rush,Apathie,EditRevert},
      yticklabels={Rush (fen.),Apathie,edit$\rightarrow$revert},
      ytick=data,
      xlabel={Nombre},
      tick label style={font=\scriptsize},
      label style={font=\scriptsize},
      nodes near coords,
      nodes near coords style={font=\tiny},
      enlarge y limits=0.25,
      grid=major,
      grid style={gray!25},
    ]
      \addplot[fill=red!45] coordinates {
        (217,EditRevert)
        (250,Apathie)
        (\RushFenetres,Rush)
      };
    \end{axis}
  \end{tikzpicture}
  \caption{Alertes CEP phase~2 (edit$\rightarrow$revert, apathie, rush).}
  \label{fig:alertes}
\end{figure}}

% --- fig08 : rush collectif (top heures) -------------------------------------
\newcommand{\figHuitRush}{%
\begin{figure}[!t]
  \centering
  \begin{tikzpicture}
    \begin{axis}[
      width=\linewidth,
      height=4.2cm,
      ybar,
      bar width=10pt,
      ymin=0,
      ymax=70,
      ylabel={Soumetteurs distincts},
      xlabel={Heure (27/04)},
      tick label style={font=\tiny, rotate=30, anchor=east},
      label style={font=\scriptsize},
      symbolic x coords={11h,13h,14h,15h,16h,17h,18h,20h},
      xtick=data,
      grid=major,
      grid style={gray!25},
      nodes near coords,
      nodes near coords style={font=\tiny},
      enlarge x limits=0.08,
    ]
      \addplot[fill=orange!70] coordinates {
        (11h,42) (13h,60) (14h,51) (15h,59) (16h,48) (17h,53) (18h,44) (20h,44)
      };
      \node[font=\tiny] at (axis cs:13h,60) [above] {pic 60};
    \end{axis}
  \end{tikzpicture}
  \caption{Rush collectif~: soumetteurs distincts par heure (pic 60, 13~h).}
  \label{fig:rush}
\end{figure}}
 puis \figXX... dans le document
% Données extraites du pipeline sur log.txt (run juin 2026) :
% counters, gauges, histogrammes, summaries, alertes JSON.
\newcommand{\RushFenetres}{24} % len(rush_collectif) dans perforce_events_alertes.json

% Coordonnées sûres pour pgfplots (éviter accents / math dans symbolic coords)
\pgfplotsset{
  /pgf/number format/.cd,
  set thousands separator={\,},
}

% --- fig00 : architecture pipeline -------------------------------------------
\newcommand{\figZeroArchitecture}{%
\begin{figure}[!t]
  \centering
  \resizebox{\linewidth}{!}{%
  \begin{tikzpicture}[
    font=\scriptsize,
    box/.style={draw, rounded corners=2pt, align=center, inner sep=4pt,
                minimum width=3.4cm, fill=white},
    small/.style={box, minimum width=3.0cm, font=\tiny},
    arr/.style={-{Stealth[length=2mm]}, thick, blue!60!black},
    darr/.style={-{Stealth[length=2mm]}, thick, gray!55},
    lbl/.style={font=\tiny, text=gray}
  ]
    \node[box, fill=blue!8, minimum width=4.2cm] (log) {\textbf{log.txt}\\{\tiny entrée (${\sim}$1,7~Go)}};
    \node[box, below=0.45cm of log] (parse) {\textbf{LecteurLog}\\{\tiny bloc = cmd + completed + \texttt{---} (pid)}};
    \node[box, below=0.45cm of parse] (dedup) {\textbf{Déduplication}\\{\tiny timestamp|pid|cmd|args|fichiers}};
    \node[small, below left=0.75cm and 0.15cm of dedup, fill=orange!10] (ds)
      {\textbf{Export v3}\\{\tiny JSON / CSV (27 col.)}};
    \node[box, below right=0.55cm and -0.1cm of dedup, fill=blue!6,
          minimum width=3.6cm] (bb)
      {\textbf{BeepBeep 3.13}\\{\tiny \texttt{QueueSource} $\rightarrow$ processeurs}};
    \node[small, below left=0.55cm and 0.05cm of bb, fill=green!8] (p1)
      {\textbf{Phase 1}\\{\tiny Counter, Gauge,\\Histogram, Summary}};
    \node[small, below right=0.55cm and 0.05cm of bb, fill=red!6] (p2)
      {\textbf{Phase 2 (CEP)}\\{\tiny \texttt{DetecteurPatterns}\\edit$\rightarrow$revert · apathie · rush}};
    \draw[arr] (log) -- (parse);
    \draw[arr] (parse) -- (dedup);
    \draw[darr] (dedup.south) -- ++(0,-0.22) -| (ds.north);
    \draw[arr] (dedup.south) -- ++(0,-0.22) -| (bb.north);
    \draw[arr] (bb.south) -- ++(0,-0.18) -| (p1.north);
    \draw[arr] (bb.south) -- ++(0,-0.18) -| (p2.north);
    \node[lbl, below=0.12cm of ds] {\texttt{perforce\_events\_v3}};
    \node[lbl, below=0.12cm of p1] {4 JSON OpenMetrics};
    \node[lbl, below=0.12cm of p2] {\texttt{perforce\_events\_alertes.json}};
  \end{tikzpicture}%
  }
  \caption{Pipeline en flux~: parse $\rightarrow$ dédup $\rightarrow$ export v3 et analyse \beepbeep{} (phases~1 et~2).}
  \label{fig:arch}
\end{figure}}

% --- fig01 : compteurs -------------------------------------------------------
\newcommand{\figUnCompteurs}{%
\begin{figure}[!t]
  \centering
  \begin{tikzpicture}
    \begin{axis}[
      width=\linewidth,
      height=4.2cm,
      ybar,
      bar width=14pt,
      ymin=0,
      symbolic x coords={Evts,Edits,Submits},
      xticklabels={Événements,Edits,Submits},
      xtick=data,
      ylabel={Nombre},
      tick label style={font=\scriptsize},
      label style={font=\scriptsize},
      nodes near coords,
      nodes near coords style={font=\tiny},
      enlarge x limits=0.2,
      grid=major,
      grid style={gray!25},
    ]
      \addplot[fill=blue!55!black] coordinates {
        (Evts,16262) (Edits,4069) (Submits,1224)
      };
    \end{axis}
  \end{tikzpicture}
  \caption{Compteurs globaux (phase~1).}
  \label{fig:compteurs}
\end{figure}}

% --- fig02 : submits par heure (points clés 27/04) ---------------------------
\newcommand{\figDeuxSubmitsHeure}{%
\begin{figure}[!t]
  \centering
  \begin{tikzpicture}
    \begin{axis}[
      width=\linewidth,
      height=4cm,
      xlabel={Heure (27/04)},
      ylabel={Submits},
      tick label style={font=\scriptsize},
      label style={font=\scriptsize},
      grid=major,
      grid style={gray!25},
      xmin=9, xmax=21, ymin=0, ymax=120,
      xtick={9,10,11,12,13,14,15,16,17,18,19,20,21},
      xticklabels={9h,10h,11h,12h,13h,14h,\textbf{15h},16h,17h,18h,19h,20h,21h},
    ]
      \addplot[thick, blue!70!black, mark=*] coordinates {
        (9,25) (10,48) (11,86) (12,63) (13,92) (14,95) (15,107)
        (16,81) (17,96) (18,75) (19,75) (20,81) (21,60)
      };
      \node[font=\tiny, anchor=south west] at (axis cs:15,107) {pic 107};
    \end{axis}
  \end{tikzpicture}
  \caption{Soumissions par heure (pic 107 à 15~h).}
  \label{fig:submits_heure}
\end{figure}}

% --- fig03 : jauge fichiers / utilisateurs ouverts ---------------------------
\newcommand{\figTroisGaugeOuverts}{%
\begin{figure}[!t]
  \centering
  \begin{tikzpicture}
    \begin{axis}[
      width=\linewidth,
      height=4cm,
      xlabel={Heure (27/04)},
      ylabel={Pic horaire},
      tick label style={font=\scriptsize},
      label style={font=\scriptsize},
      legend style={font=\tiny, at={(0.02,0.98)}, anchor=north west},
      grid=major,
      grid style={gray!25},
      xmin=9, xmax=20, ymin=0, ymax=520,
      xtick={10,11,12,13,14,15,16,17,18,19},
      xticklabels={10h,11h,12h,13h,14h,\textbf{15h},16h,17h,18h,19h},
    ]
      \addplot[thick, green!50!black, mark=square*] coordinates {
        (10,245) (11,311) (12,347) (13,334) (14,343) (15,294)
        (16,500) (17,286) (18,263) (19,317)
      };
      \addlegendentry{Fichiers ouverts}
      \addplot[thick, orange!80!black, mark=triangle*] coordinates {
        (10,53) (11,54) (12,64) (13,74) (14,79) (15,78)
        (16,75) (17,71) (18,72) (19,74)
      };
      \addlegendentry{Utilisateurs}
      \node[font=\tiny] at (axis cs:16,500) [above] {500};
      \node[font=\tiny] at (axis cs:14,79) [above] {79};
    \end{axis}
  \end{tikzpicture}
  \caption{Jauge — fichiers et utilisateurs en edit (pics horaires).}
  \label{fig:gauge_ouverts}
\end{figure}}

% --- fig02+03 combinées (fig:activite) ---------------------------------------
\newcommand{\figActivite}{%
\begin{figure}[!t]
  \centering
  \begin{subfigure}{\linewidth}
    \centering
    \begin{tikzpicture}
      \begin{axis}[
        width=\linewidth, height=3.2cm,
        xlabel={Heure (27/04)}, ylabel={Submits},
        tick label style={font=\tiny}, label style={font=\scriptsize},
        grid=major, grid style={gray!25},
        xmin=9, xmax=21, ymin=0, ymax=120,
        xtick={9,10,11,12,13,14,15,16,17,18,19,20,21},
        xticklabels={9h,10h,11h,12h,13h,14h,\textbf{15h},16h,17h,18h,19h,20h,21h},
      ]
        \addplot[thick, blue!70!black, mark=*] coordinates {
          (9,25) (10,48) (11,86) (12,63) (13,92) (14,95) (15,107)
          (16,81) (17,96) (18,75) (19,75) (20,81) (21,60)
        };
      \end{axis}
    \end{tikzpicture}
    \caption{Soumissions/heure (pic 107, 15~h).}
    \label{fig:submits_heure}
  \end{subfigure}\\[0.4em]
  \begin{subfigure}{\linewidth}
    \centering
    \begin{tikzpicture}
      \begin{axis}[
        width=\linewidth, height=3.2cm,
        xlabel={Heure (27/04)}, ylabel={Fichiers},
        tick label style={font=\tiny}, label style={font=\scriptsize},
        grid=major, grid style={gray!25},
        xmin=9, xmax=20, ymin=0, ymax=520,
        xtick={10,11,12,13,14,15,16,17,18,19},
        xticklabels={10h,11h,12h,13h,14h,15h,\textbf{16h},17h,18h,19h},
      ]
        \addplot[thick, green!50!black, mark=square*] coordinates {
          (10,245) (11,311) (12,347) (13,334) (14,343) (15,294)
          (16,500) (17,286) (18,263) (19,317)
        };
      \end{axis}
    \end{tikzpicture}
    \caption{Fichiers ouverts (pic 500).}
    \label{fig:gauge_ouverts}
  \end{subfigure}
  \caption{Activité sur 24~h autour de l'échéance.}
  \label{fig:activite}
\end{figure}}

% --- fig04 : histogramme durées (16 232 obs., deltas JSON) --------------------
\newcommand{\figQuatreHistogramme}{%
\begin{figure}[!t]
  \centering
  \begin{tikzpicture}
    \begin{axis}[
      width=\linewidth,
      height=4.5cm,
      ybar,
      bar width=5pt,
      ymin=0,
      ymax=6500,
      ylabel={Obs.},
      xlabel={Bucket (s)},
      tick label style={font=\tiny, rotate=45, anchor=east},
      label style={font=\scriptsize},
      symbolic x coords={%
        le005,le01,le025,le05,le1,le2,le5,le10,le30,le60,le120,le300,le600,le1200},
      xticklabels={.05,.1,.25,.5,1,2,5,10,30,60,120,300,600,1200},
      xtick=data,
      enlarge x limits=0.02,
      grid=major,
      grid style={gray!25},
      nodes near coords={},
    ]
      % deltas exacts (cumulatif → par bucket), total 16 232
      \addplot[fill=blue!50] coordinates {
        (le005,2180) (le01,6114) (le025,6202) (le05,584) (le1,301) (le2,160) (le5,189)
        (le10,87) (le30,136) (le60,81) (le120,65) (le300,48) (le600,38) (le1200,27)
      };
    \end{axis}
  \end{tikzpicture}
  \caption{Histogramme des durées (16\,232 obs., 15 buckets).}
  \label{fig:durees}
\end{figure}}

% --- fig05 : quantiles par commande ------------------------------------------
\newcommand{\figCinqQuantiles}{%
\begin{figure}[!t]
  \centering
  \begin{tikzpicture}
    \begin{axis}[
      width=\linewidth,
      height=4.2cm,
      ybar,
      bar width=4pt,
      ymin=0,
      ymax=42,
      ylabel={Secondes},
      symbolic x coords={edit,submit,add,resolve},
      xtick=data,
      legend style={font=\tiny, at={(0.98,0.98)}, anchor=north east},
      tick label style={font=\scriptsize},
      label style={font=\scriptsize},
      grid=major,
      grid style={gray!25},
      enlarge x limits=0.15,
    ]
      \addplot[fill=blue!40] coordinates {(edit,0.110) (submit,0.125) (add,0.125) (resolve,0.078)};
      \addlegendentry{p50}
      \addplot[fill=orange!60] coordinates {(edit,0.204) (submit,0.250) (add,0.452) (resolve,0.094)};
      \addlegendentry{p95}
      \addplot[fill=red!50] coordinates {(edit,0.698) (submit,0.821) (add,39.349) (resolve,0.172)};
      \addlegendentry{p99}
    \end{axis}
  \end{tikzpicture}
  \caption{Quantiles Summary par commande (p50, p95, p99).}
  \label{fig:quantiles}
\end{figure}}

% --- fig06 : zoom sync (échelle log) -----------------------------------------
\newcommand{\figSixQuantilesSync}{%
\begin{figure}[!t]
  \centering
  \begin{tikzpicture}
    \begin{axis}[
      width=0.85\linewidth,
      height=4cm,
      ybar,
      bar width=18pt,
      ymin=0.1,
      ymax=2000,
      ymode=log,
      log basis y=10,
      ylabel={Secondes (log)},
      symbolic x coords={p50,p95,p99},
      xtick=data,
      tick label style={font=\scriptsize},
      label style={font=\scriptsize},
      grid=major,
      grid style={gray!25},
      nodes near coords={\pgfmathprintnumber\pgfplotspointmeta},
      nodes near coords style={font=\tiny, /pgf/number format/fixed},
      enlarge x limits=0.3,
    ]
      \addplot[fill=purple!45] coordinates {(p50,0.188) (p95,104.65) (p99,1101.74)};
    \end{axis}
  \end{tikzpicture}
  \caption{Quantiles \texttt{user-sync} (échelle logarithmique).}
  \label{fig:quantiles_sync}
\end{figure}}

% --- fig05+06 combinées ------------------------------------------------------
\newcommand{\figQuantilesCombine}{%
\begin{figure}[!t]
  \centering
  \begin{subfigure}{0.48\linewidth}
    \centering
    \resizebox{\linewidth}{!}{%
    \begin{tikzpicture}
      \begin{axis}[
        width=6cm, height=4cm, ybar, bar width=4pt,
        ymin=0, ymax=42, ylabel={s},
        symbolic x coords={edit,submit,add,res},
        xtick=data,
        tick label style={font=\tiny},
        legend style={font=\tiny, legend columns=3, at={(0.5,1.12)}, anchor=south},
        enlarge x limits=0.12,
      ]
        \addplot[fill=blue!40] coordinates {(edit,0.110) (submit,0.125) (add,0.125) (res,0.078)};
        \addlegendentry{p50}
        \addplot[fill=orange!60] coordinates {(edit,0.204) (submit,0.250) (add,0.452) (res,0.094)};
        \addlegendentry{p95}
        \addplot[fill=red!50] coordinates {(edit,0.698) (submit,0.821) (add,39.349) (res,0.172)};
        \addlegendentry{p99}
      \end{axis}
    \end{tikzpicture}}
    \caption{Par commande.}
    \label{fig:quantiles}
  \end{subfigure}\hfill
  \begin{subfigure}{0.48\linewidth}
    \centering
    \resizebox{\linewidth}{!}{%
    \begin{tikzpicture}
      \begin{axis}[
        width=5cm, height=4cm, ybar, bar width=12pt,
        ymin=0.1, ymax=2000, ymode=log, log basis y=10,
        ylabel={s (log)},
        symbolic x coords={p50,p95,p99},
        xtick=data,
        tick label style={font=\tiny},
        nodes near coords={\pgfmathprintnumber\pgfplotspointmeta},
        nodes near coords style={font=\tiny},
        enlarge x limits=0.25,
      ]
        \addplot[fill=purple!45] coordinates {(p50,0.188) (p95,104.65) (p99,1101.74)};
      \end{axis}
    \end{tikzpicture}}
    \caption{Zoom \texttt{sync}.}
    \label{fig:quantiles_sync}
  \end{subfigure}
  \caption{Summary phase~1 (p50, p95, p99).}
\end{figure}}

% --- fig07 : alertes phase 2 (diagramme en barres, plus lisible que tarte) ---
\newcommand{\figSeptAlertes}{%
\begin{figure}[!t]
  \centering
  \begin{tikzpicture}
    \begin{axis}[
      width=\linewidth,
      height=3.8cm,
      xbar,
      bar width=10pt,
      xmin=0,
      symbolic y coords={Rush,Apathie,EditRevert},
      yticklabels={Rush (fen.),Apathie,edit$\rightarrow$revert},
      ytick=data,
      xlabel={Nombre},
      tick label style={font=\scriptsize},
      label style={font=\scriptsize},
      nodes near coords,
      nodes near coords style={font=\tiny},
      enlarge y limits=0.25,
      grid=major,
      grid style={gray!25},
    ]
      \addplot[fill=red!45] coordinates {
        (217,EditRevert)
        (250,Apathie)
        (\RushFenetres,Rush)
      };
    \end{axis}
  \end{tikzpicture}
  \caption{Alertes CEP phase~2 (edit$\rightarrow$revert, apathie, rush).}
  \label{fig:alertes}
\end{figure}}

% --- fig08 : rush collectif (top heures) -------------------------------------
\newcommand{\figHuitRush}{%
\begin{figure}[!t]
  \centering
  \begin{tikzpicture}
    \begin{axis}[
      width=\linewidth,
      height=4.2cm,
      ybar,
      bar width=10pt,
      ymin=0,
      ymax=70,
      ylabel={Soumetteurs distincts},
      xlabel={Heure (27/04)},
      tick label style={font=\tiny, rotate=30, anchor=east},
      label style={font=\scriptsize},
      symbolic x coords={11h,13h,14h,15h,16h,17h,18h,20h},
      xtick=data,
      grid=major,
      grid style={gray!25},
      nodes near coords,
      nodes near coords style={font=\tiny},
      enlarge x limits=0.08,
    ]
      \addplot[fill=orange!70] coordinates {
        (11h,42) (13h,60) (14h,51) (15h,59) (16h,48) (17h,53) (18h,44) (20h,44)
      };
      \node[font=\tiny] at (axis cs:13h,60) [above] {pic 60};
    \end{axis}
  \end{tikzpicture}
  \caption{Rush collectif~: soumetteurs distincts par heure (pic 60, 13~h).}
  \label{fig:rush}
\end{figure}}
 puis \figXX... dans le document
% Données extraites du pipeline sur log.txt (run juin 2026) :
% counters, gauges, histogrammes, summaries, alertes JSON.
\newcommand{\RushFenetres}{24} % len(rush_collectif) dans perforce_events_alertes.json

% Coordonnées sûres pour pgfplots (éviter accents / math dans symbolic coords)
\pgfplotsset{
  /pgf/number format/.cd,
  set thousands separator={\,},
}

% --- fig00 : architecture pipeline -------------------------------------------
\newcommand{\figZeroArchitecture}{%
\begin{figure}[!t]
  \centering
  \resizebox{\linewidth}{!}{%
  \begin{tikzpicture}[
    font=\scriptsize,
    box/.style={draw, rounded corners=2pt, align=center, inner sep=4pt,
                minimum width=3.4cm, fill=white},
    small/.style={box, minimum width=3.0cm, font=\tiny},
    arr/.style={-{Stealth[length=2mm]}, thick, blue!60!black},
    darr/.style={-{Stealth[length=2mm]}, thick, gray!55},
    lbl/.style={font=\tiny, text=gray}
  ]
    \node[box, fill=blue!8, minimum width=4.2cm] (log) {\textbf{log.txt}\\{\tiny entrée (${\sim}$1,7~Go)}};
    \node[box, below=0.45cm of log] (parse) {\textbf{LecteurLog}\\{\tiny bloc = cmd + completed + \texttt{---} (pid)}};
    \node[box, below=0.45cm of parse] (dedup) {\textbf{Déduplication}\\{\tiny timestamp|pid|cmd|args|fichiers}};
    \node[small, below left=0.75cm and 0.15cm of dedup, fill=orange!10] (ds)
      {\textbf{Export v3}\\{\tiny JSON / CSV (27 col.)}};
    \node[box, below right=0.55cm and -0.1cm of dedup, fill=blue!6,
          minimum width=3.6cm] (bb)
      {\textbf{BeepBeep 3.13}\\{\tiny \texttt{QueueSource} $\rightarrow$ processeurs}};
    \node[small, below left=0.55cm and 0.05cm of bb, fill=green!8] (p1)
      {\textbf{Phase 1}\\{\tiny Counter, Gauge,\\Histogram, Summary}};
    \node[small, below right=0.55cm and 0.05cm of bb, fill=red!6] (p2)
      {\textbf{Phase 2 (CEP)}\\{\tiny \texttt{DetecteurPatterns}\\edit$\rightarrow$revert · apathie · rush}};
    \draw[arr] (log) -- (parse);
    \draw[arr] (parse) -- (dedup);
    \draw[darr] (dedup.south) -- ++(0,-0.22) -| (ds.north);
    \draw[arr] (dedup.south) -- ++(0,-0.22) -| (bb.north);
    \draw[arr] (bb.south) -- ++(0,-0.18) -| (p1.north);
    \draw[arr] (bb.south) -- ++(0,-0.18) -| (p2.north);
    \node[lbl, below=0.12cm of ds] {\texttt{perforce\_events\_v3}};
    \node[lbl, below=0.12cm of p1] {4 JSON OpenMetrics};
    \node[lbl, below=0.12cm of p2] {\texttt{perforce\_events\_alertes.json}};
  \end{tikzpicture}%
  }
  \caption{Pipeline en flux~: parse $\rightarrow$ dédup $\rightarrow$ export v3 et analyse \beepbeep{} (phases~1 et~2).}
  \label{fig:arch}
\end{figure}}

% --- fig01 : compteurs -------------------------------------------------------
\newcommand{\figUnCompteurs}{%
\begin{figure}[!t]
  \centering
  \begin{tikzpicture}
    \begin{axis}[
      width=\linewidth,
      height=4.2cm,
      ybar,
      bar width=14pt,
      ymin=0,
      symbolic x coords={Evts,Edits,Submits},
      xticklabels={Événements,Edits,Submits},
      xtick=data,
      ylabel={Nombre},
      tick label style={font=\scriptsize},
      label style={font=\scriptsize},
      nodes near coords,
      nodes near coords style={font=\tiny},
      enlarge x limits=0.2,
      grid=major,
      grid style={gray!25},
    ]
      \addplot[fill=blue!55!black] coordinates {
        (Evts,16262) (Edits,4069) (Submits,1224)
      };
    \end{axis}
  \end{tikzpicture}
  \caption{Compteurs globaux (phase~1).}
  \label{fig:compteurs}
\end{figure}}

% --- fig02 : submits par heure (points clés 27/04) ---------------------------
\newcommand{\figDeuxSubmitsHeure}{%
\begin{figure}[!t]
  \centering
  \begin{tikzpicture}
    \begin{axis}[
      width=\linewidth,
      height=4cm,
      xlabel={Heure (27/04)},
      ylabel={Submits},
      tick label style={font=\scriptsize},
      label style={font=\scriptsize},
      grid=major,
      grid style={gray!25},
      xmin=9, xmax=21, ymin=0, ymax=120,
      xtick={9,10,11,12,13,14,15,16,17,18,19,20,21},
      xticklabels={9h,10h,11h,12h,13h,14h,\textbf{15h},16h,17h,18h,19h,20h,21h},
    ]
      \addplot[thick, blue!70!black, mark=*] coordinates {
        (9,25) (10,48) (11,86) (12,63) (13,92) (14,95) (15,107)
        (16,81) (17,96) (18,75) (19,75) (20,81) (21,60)
      };
      \node[font=\tiny, anchor=south west] at (axis cs:15,107) {pic 107};
    \end{axis}
  \end{tikzpicture}
  \caption{Soumissions par heure (pic 107 à 15~h).}
  \label{fig:submits_heure}
\end{figure}}

% --- fig03 : jauge fichiers / utilisateurs ouverts ---------------------------
\newcommand{\figTroisGaugeOuverts}{%
\begin{figure}[!t]
  \centering
  \begin{tikzpicture}
    \begin{axis}[
      width=\linewidth,
      height=4cm,
      xlabel={Heure (27/04)},
      ylabel={Pic horaire},
      tick label style={font=\scriptsize},
      label style={font=\scriptsize},
      legend style={font=\tiny, at={(0.02,0.98)}, anchor=north west},
      grid=major,
      grid style={gray!25},
      xmin=9, xmax=20, ymin=0, ymax=520,
      xtick={10,11,12,13,14,15,16,17,18,19},
      xticklabels={10h,11h,12h,13h,14h,\textbf{15h},16h,17h,18h,19h},
    ]
      \addplot[thick, green!50!black, mark=square*] coordinates {
        (10,245) (11,311) (12,347) (13,334) (14,343) (15,294)
        (16,500) (17,286) (18,263) (19,317)
      };
      \addlegendentry{Fichiers ouverts}
      \addplot[thick, orange!80!black, mark=triangle*] coordinates {
        (10,53) (11,54) (12,64) (13,74) (14,79) (15,78)
        (16,75) (17,71) (18,72) (19,74)
      };
      \addlegendentry{Utilisateurs}
      \node[font=\tiny] at (axis cs:16,500) [above] {500};
      \node[font=\tiny] at (axis cs:14,79) [above] {79};
    \end{axis}
  \end{tikzpicture}
  \caption{Jauge — fichiers et utilisateurs en edit (pics horaires).}
  \label{fig:gauge_ouverts}
\end{figure}}

% --- fig02+03 combinées (fig:activite) ---------------------------------------
\newcommand{\figActivite}{%
\begin{figure}[!t]
  \centering
  \begin{subfigure}{\linewidth}
    \centering
    \begin{tikzpicture}
      \begin{axis}[
        width=\linewidth, height=3.2cm,
        xlabel={Heure (27/04)}, ylabel={Submits},
        tick label style={font=\tiny}, label style={font=\scriptsize},
        grid=major, grid style={gray!25},
        xmin=9, xmax=21, ymin=0, ymax=120,
        xtick={9,10,11,12,13,14,15,16,17,18,19,20,21},
        xticklabels={9h,10h,11h,12h,13h,14h,\textbf{15h},16h,17h,18h,19h,20h,21h},
      ]
        \addplot[thick, blue!70!black, mark=*] coordinates {
          (9,25) (10,48) (11,86) (12,63) (13,92) (14,95) (15,107)
          (16,81) (17,96) (18,75) (19,75) (20,81) (21,60)
        };
      \end{axis}
    \end{tikzpicture}
    \caption{Soumissions/heure (pic 107, 15~h).}
    \label{fig:submits_heure}
  \end{subfigure}\\[0.4em]
  \begin{subfigure}{\linewidth}
    \centering
    \begin{tikzpicture}
      \begin{axis}[
        width=\linewidth, height=3.2cm,
        xlabel={Heure (27/04)}, ylabel={Fichiers},
        tick label style={font=\tiny}, label style={font=\scriptsize},
        grid=major, grid style={gray!25},
        xmin=9, xmax=20, ymin=0, ymax=520,
        xtick={10,11,12,13,14,15,16,17,18,19},
        xticklabels={10h,11h,12h,13h,14h,15h,\textbf{16h},17h,18h,19h},
      ]
        \addplot[thick, green!50!black, mark=square*] coordinates {
          (10,245) (11,311) (12,347) (13,334) (14,343) (15,294)
          (16,500) (17,286) (18,263) (19,317)
        };
      \end{axis}
    \end{tikzpicture}
    \caption{Fichiers ouverts (pic 500).}
    \label{fig:gauge_ouverts}
  \end{subfigure}
  \caption{Activité sur 24~h autour de l'échéance.}
  \label{fig:activite}
\end{figure}}

% --- fig04 : histogramme durées (16 232 obs., deltas JSON) --------------------
\newcommand{\figQuatreHistogramme}{%
\begin{figure}[!t]
  \centering
  \begin{tikzpicture}
    \begin{axis}[
      width=\linewidth,
      height=4.5cm,
      ybar,
      bar width=5pt,
      ymin=0,
      ymax=6500,
      ylabel={Obs.},
      xlabel={Bucket (s)},
      tick label style={font=\tiny, rotate=45, anchor=east},
      label style={font=\scriptsize},
      symbolic x coords={%
        le005,le01,le025,le05,le1,le2,le5,le10,le30,le60,le120,le300,le600,le1200},
      xticklabels={.05,.1,.25,.5,1,2,5,10,30,60,120,300,600,1200},
      xtick=data,
      enlarge x limits=0.02,
      grid=major,
      grid style={gray!25},
      nodes near coords={},
    ]
      % deltas exacts (cumulatif → par bucket), total 16 232
      \addplot[fill=blue!50] coordinates {
        (le005,2180) (le01,6114) (le025,6202) (le05,584) (le1,301) (le2,160) (le5,189)
        (le10,87) (le30,136) (le60,81) (le120,65) (le300,48) (le600,38) (le1200,27)
      };
    \end{axis}
  \end{tikzpicture}
  \caption{Histogramme des durées (16\,232 obs., 15 buckets).}
  \label{fig:durees}
\end{figure}}

% --- fig05 : quantiles par commande ------------------------------------------
\newcommand{\figCinqQuantiles}{%
\begin{figure}[!t]
  \centering
  \begin{tikzpicture}
    \begin{axis}[
      width=\linewidth,
      height=4.2cm,
      ybar,
      bar width=4pt,
      ymin=0,
      ymax=42,
      ylabel={Secondes},
      symbolic x coords={edit,submit,add,resolve},
      xtick=data,
      legend style={font=\tiny, at={(0.98,0.98)}, anchor=north east},
      tick label style={font=\scriptsize},
      label style={font=\scriptsize},
      grid=major,
      grid style={gray!25},
      enlarge x limits=0.15,
    ]
      \addplot[fill=blue!40] coordinates {(edit,0.110) (submit,0.125) (add,0.125) (resolve,0.078)};
      \addlegendentry{p50}
      \addplot[fill=orange!60] coordinates {(edit,0.204) (submit,0.250) (add,0.452) (resolve,0.094)};
      \addlegendentry{p95}
      \addplot[fill=red!50] coordinates {(edit,0.698) (submit,0.821) (add,39.349) (resolve,0.172)};
      \addlegendentry{p99}
    \end{axis}
  \end{tikzpicture}
  \caption{Quantiles Summary par commande (p50, p95, p99).}
  \label{fig:quantiles}
\end{figure}}

% --- fig06 : zoom sync (échelle log) -----------------------------------------
\newcommand{\figSixQuantilesSync}{%
\begin{figure}[!t]
  \centering
  \begin{tikzpicture}
    \begin{axis}[
      width=0.85\linewidth,
      height=4cm,
      ybar,
      bar width=18pt,
      ymin=0.1,
      ymax=2000,
      ymode=log,
      log basis y=10,
      ylabel={Secondes (log)},
      symbolic x coords={p50,p95,p99},
      xtick=data,
      tick label style={font=\scriptsize},
      label style={font=\scriptsize},
      grid=major,
      grid style={gray!25},
      nodes near coords={\pgfmathprintnumber\pgfplotspointmeta},
      nodes near coords style={font=\tiny, /pgf/number format/fixed},
      enlarge x limits=0.3,
    ]
      \addplot[fill=purple!45] coordinates {(p50,0.188) (p95,104.65) (p99,1101.74)};
    \end{axis}
  \end{tikzpicture}
  \caption{Quantiles \texttt{user-sync} (échelle logarithmique).}
  \label{fig:quantiles_sync}
\end{figure}}

% --- fig05+06 combinées ------------------------------------------------------
\newcommand{\figQuantilesCombine}{%
\begin{figure}[!t]
  \centering
  \begin{subfigure}{0.48\linewidth}
    \centering
    \resizebox{\linewidth}{!}{%
    \begin{tikzpicture}
      \begin{axis}[
        width=6cm, height=4cm, ybar, bar width=4pt,
        ymin=0, ymax=42, ylabel={s},
        symbolic x coords={edit,submit,add,res},
        xtick=data,
        tick label style={font=\tiny},
        legend style={font=\tiny, legend columns=3, at={(0.5,1.12)}, anchor=south},
        enlarge x limits=0.12,
      ]
        \addplot[fill=blue!40] coordinates {(edit,0.110) (submit,0.125) (add,0.125) (res,0.078)};
        \addlegendentry{p50}
        \addplot[fill=orange!60] coordinates {(edit,0.204) (submit,0.250) (add,0.452) (res,0.094)};
        \addlegendentry{p95}
        \addplot[fill=red!50] coordinates {(edit,0.698) (submit,0.821) (add,39.349) (res,0.172)};
        \addlegendentry{p99}
      \end{axis}
    \end{tikzpicture}}
    \caption{Par commande.}
    \label{fig:quantiles}
  \end{subfigure}\hfill
  \begin{subfigure}{0.48\linewidth}
    \centering
    \resizebox{\linewidth}{!}{%
    \begin{tikzpicture}
      \begin{axis}[
        width=5cm, height=4cm, ybar, bar width=12pt,
        ymin=0.1, ymax=2000, ymode=log, log basis y=10,
        ylabel={s (log)},
        symbolic x coords={p50,p95,p99},
        xtick=data,
        tick label style={font=\tiny},
        nodes near coords={\pgfmathprintnumber\pgfplotspointmeta},
        nodes near coords style={font=\tiny},
        enlarge x limits=0.25,
      ]
        \addplot[fill=purple!45] coordinates {(p50,0.188) (p95,104.65) (p99,1101.74)};
      \end{axis}
    \end{tikzpicture}}
    \caption{Zoom \texttt{sync}.}
    \label{fig:quantiles_sync}
  \end{subfigure}
  \caption{Summary phase~1 (p50, p95, p99).}
\end{figure}}

% --- fig07 : alertes phase 2 (diagramme en barres, plus lisible que tarte) ---
\newcommand{\figSeptAlertes}{%
\begin{figure}[!t]
  \centering
  \begin{tikzpicture}
    \begin{axis}[
      width=\linewidth,
      height=3.8cm,
      xbar,
      bar width=10pt,
      xmin=0,
      symbolic y coords={Rush,Apathie,EditRevert},
      yticklabels={Rush (fen.),Apathie,edit$\rightarrow$revert},
      ytick=data,
      xlabel={Nombre},
      tick label style={font=\scriptsize},
      label style={font=\scriptsize},
      nodes near coords,
      nodes near coords style={font=\tiny},
      enlarge y limits=0.25,
      grid=major,
      grid style={gray!25},
    ]
      \addplot[fill=red!45] coordinates {
        (217,EditRevert)
        (250,Apathie)
        (\RushFenetres,Rush)
      };
    \end{axis}
  \end{tikzpicture}
  \caption{Alertes CEP phase~2 (edit$\rightarrow$revert, apathie, rush).}
  \label{fig:alertes}
\end{figure}}

% --- fig08 : rush collectif (top heures) -------------------------------------
\newcommand{\figHuitRush}{%
\begin{figure}[!t]
  \centering
  \begin{tikzpicture}
    \begin{axis}[
      width=\linewidth,
      height=4.2cm,
      ybar,
      bar width=10pt,
      ymin=0,
      ymax=70,
      ylabel={Soumetteurs distincts},
      xlabel={Heure (27/04)},
      tick label style={font=\tiny, rotate=30, anchor=east},
      label style={font=\scriptsize},
      symbolic x coords={11h,13h,14h,15h,16h,17h,18h,20h},
      xtick=data,
      grid=major,
      grid style={gray!25},
      nodes near coords,
      nodes near coords style={font=\tiny},
      enlarge x limits=0.08,
    ]
      \addplot[fill=orange!70] coordinates {
        (11h,42) (13h,60) (14h,51) (15h,59) (16h,48) (17h,53) (18h,44) (20h,44)
      };
      \node[font=\tiny] at (axis cs:13h,60) [above] {pic 60};
    \end{axis}
  \end{tikzpicture}
  \caption{Rush collectif~: soumetteurs distincts par heure (pic 60, 13~h).}
  \label{fig:rush}
\end{figure}}


\begin{document}

\title{Detecting Collaboration Antipatterns in Perforce Server Logs\\
with Stream Metrics and Complex Event Processing}

\author{%
\IEEEauthorblockN{Hanifah ALPHA-BODA}
\IEEEauthorblockA{\textit{Laboratoire d'informatique formelle}\\
\textit{Universit\'e du Qu\'ebec \`a Chicoutimi}\\
Chicoutimi, Canada\\
haboda@etu.uqac.ca}
\and
\IEEEauthorblockN{Yannick FRANCILETTE}
\IEEEauthorblockA{\textit{Multi-user Interaction Experience Lab}\\
\textit{Universit\'e du Qu\'ebec \`a Chicoutimi}\\
Chicoutimi, Canada\\
yannick\_francillette@uqac.ca}
\and
\IEEEauthorblockN{Sylvain HALL\'E}
\IEEEauthorblockA{\textit{Laboratoire d'informatique formelle}\\
\textit{Universit\'e du Qu\'ebec \`a Chicoutimi}\\
Chicoutimi, Canada\\
shalle@uqac.ca}
}

\maketitle

\begin{abstract}
Les serveurs de gestion de versions enregistrent chaque action d'une
équipe, mais ces journaux sont rarement exploités pour analyser la
collaboration. Dans l'industrie du jeu vidéo, des équipes nombreuses
manipulent des \textbf{fichiers binaires} non fusionnables au sein de
dépôts volumineux~: les signaux utiles proviennent des séquences
d'opérations (modification, synchronisation, soumission, annulation) et
de leurs durées, non de comparaisons de contenu textuel.

Nous présentons \toolname{}, une \textbf{contribution logicielle} de
détection d'\emph{antipatterns de collaboration} à partir des journaux du
serveur Perforce. Cinq motifs du catalogue ReliSA (sur 168) sont
opérationnalisés. L'outil assemble des blocs d'événements puis enchaîne
deux traitements~: métriques agrégées inspirées de
Prometheus~\cite{brazil2018prometheus,prometheus_metric_types}, puis
traitement d'événements complexes sensible à l'ordre des
actions~\cite{halle2018beepbeep,luckham2002power}.

Sur un journal de 1,7~Go, il extrait 16\,262 événements uniques en
${\sim}$14~s, signale 217 séquences \emph{modification-puis-annulation}
et 250 fichiers abandonnés~; une jauge indépendante confirme ce
décompte (250\,=\,250). Ces résultats montrent la faisabilité d'une
analyse automatique à grande échelle et la cohérence interne du
pipeline proposé.
\end{abstract}

\begin{IEEEkeywords}
traitement d'événements complexes, fouille de dépôts logiciels,
antipatterns de collaboration, systèmes de gestion de versions,
journaux Perforce, métriques en flux
\end{IEEEkeywords}

\section{Introduction}

\textbf{Contexte.}
Le développement de jeux vidéo repose sur de grandes équipes et des
outils de gestion de versions partagés. Les retards de production ont un
coût financier majeur pour les studios~; une collaboration
dysfonctionnelle contribue souvent à ces dérapages.

\textbf{Problématique.}
La recherche s'intéresse aux \emph{antipatterns de
collaboration}~\cite{brown1998antipatterns,tamburri2019community,relisa_catalogue}~:
comportements d'équipe récurrents et contre-productifs. Les détecter tôt
permettrait une rétroaction avant qu'ils n'affectent la livraison.

\textbf{Problème.}
Les journaux du serveur Perforce (\texttt{p4d}) enregistrent chaque
\texttt{edit}, \texttt{submit}, \texttt{sync} et \texttt{revert}, mais
atteignent facilement plusieurs gigaoctets, illisibles manuellement.
Contrairement à Git, centré sur le graphe de commits, Perforce journalise
les opérations \emph{au niveau fichier}, ce qui rend observable un
fichier «~ouvert mais jamais soumis~»---inaccessible aux approches
orientées commits~\cite{hindle2008mining}. Sur des \textbf{fichiers
binaires} de moteur de jeu, l'analyse repose sur ces séquences et leurs
durées, non sur des diffs textuels.

\textbf{Question de recherche.}
Comment analyser automatiquement ces journaux pour détecter des
mauvaises pratiques de collaboration, en flux et en tenant compte de
l'ordre des événements sur chaque fichier~?

\textbf{Contributions.}
Il s'agit d'une contribution \textbf{appliquée} (outil reproductible),
non théorique. \toolname{} transforme un journal en signaux
d'antipatterns~: comportements comme un membre qui modifie beaucoup mais
partage peu, ou un fichier ouvert puis abandonné. ReliSA recense
\textbf{168} fiches~\cite{relisa_catalogue}~; nous en opérationnalisons
\textbf{cinq} (tableau~\ref{tab:relisas}), comme \textbf{proxies} dans
\texttt{p4d}. Pipeline \emph{étagé}~: métriques
Prometheus/OpenMetrics~\cite{brazil2018prometheus,openmetrics}, puis
traitement d'événements complexes avec \beepbeep{}~\cite{halle2018beepbeep}
(«~combien~?~» vs «~dans quel ordre~?~»). (1)~Blocs d'événement
(\texttt{user-*} + \texttt{completed} + \texttt{---}, \texttt{pid})~;
(2)~métriques hors ligne~; (3)~automates pour motifs temporels.

\begin{table}[!t]
  \caption{Cinq antipatterns opérationnalisés (5/168 ReliSA).}
  \label{tab:relisas}
  \centering
  \small
  \begin{tabular}{ll}
    \toprule
    \textbf{Notre détection} & \textbf{Fiche ReliSA} \\
    \midrule
    Rush de submits & Collective\_Procrastination \\
    Faible activité (Warm Bodies) & Warm\_Bodies \\
    Ratio edit/submit élevé (Lone Wolf) & Lone-Wolf \\
    Checkout sans submit (apathie) & Bystander\_Apathy \\
    Edit puis revert & Cascading\_Branches (approx.) \\
    \bottomrule
  \end{tabular}
\end{table}

\section{\'Etat de l'art}
\label{sec:etat}

\textbf{Critères de sélection.}
Travaux sur fouille de dépôts, antipatterns, observabilité par métriques
et traitement d'événements complexes, fenêtre 2018--2026 (références
fondatrices plus anciennes si
nécessaire)~\cite{luckham2002power,brown1998antipatterns}.

\textbf{Présentation.}
La fouille sur Git~\cite{hindle2008mining} analyse des commits en lot.
Les \emph{community smells}~\cite{tamburri2019community} ciblent des
signaux organisationnels souvent absents des journaux VCS. Prometheus
définit compteurs, jauges, histogrammes et
résumés~\cite{brazil2018prometheus,prometheus_metric_types}~;
OpenMetrics standardise leur
exposition~\cite{openmetrics}. Les moteurs de traitement d'événements
complexes~\cite{luckham2002power} maintiennent un état sur un flux
ordonné. \beepbeep{}~\cite{halle2018beepbeep} unifie métriques et
automates en pipelines Java composables.

\textbf{Comparaison.}
Le tableau~\ref{tab:comparaison} situe notre approche par rapport aux
travaux voisins.

\begin{table}[!t]
  \caption{Comparaison synthétique.}
  \label{tab:comparaison}
  \centering
  \scriptsize
  \setlength{\tabcolsep}{2pt}
  \begin{tabular}{@{}lccccc@{}}
    \toprule
    \textbf{Approche} & \textbf{Flot} & \textbf{Ordre} & \textbf{\texttt{p4d}} & \textbf{Go} & \textbf{Collab.} \\
    \midrule
    Fouille Git~\cite{hindle2008mining} & \texttimes & \texttimes & \texttimes & \checkmark & partiel \\
    Community smells~\cite{tamburri2019community} & \texttimes & \texttimes & \texttimes & \checkmark & \checkmark \\
    Métriques Prometheus~\cite{prometheus_metric_types} & \checkmark & \texttimes & \texttimes & \checkmark & \texttimes \\
    \textbf{\toolname{} (notre)} & \checkmark & \checkmark & \checkmark & \checkmark & \checkmark \\
    \bottomrule
  \end{tabular}
\end{table}

\textbf{Discussion.}
Notre force~: échelle (${\sim}$14~s sur 1,7~Go), granularité fichier
propre à Perforce, validation croisée jauge/automate. Compromis~: cinq
proxies sur 168 fiches ReliSA, seuils contextuels, précision des alertes
non encore mesurée.

\section{M\'ethodologie}
\label{sec:methodo}

\subsection{Mod\`ele d'\'ev\'enements et antipatterns cibl\'es}
Un événement est un tuple $e=(t,u,f,a)$ ($t$ horodatage, $u$ utilisateur,
$f$ chemin fichier, $a$ action). Pour $u$, $E(u)$ et $S(u)$ comptent
\texttt{edit} et \texttt{submit}. Les cinq antipatterns sont~:
\begin{enumerate}
\item \emph{Procrastination collective}~: taux de soumission $>\rho$ avant échéance.
\item \emph{Warm Bodies}~: $E(u)+S(u) < \theta_{\mathrm{wb}}$.
\item \emph{Lone Wolf}~: $E(u)/\max(S(u),1) \geq \tau$.
\item \emph{Bystander Apathy}~: \texttt{edit} sans \texttt{submit} ultérieur avant fin de flux.
\item \emph{Modification-puis-annulation}~: \texttt{edit} puis \texttt{revert} sans \texttt{submit}.
\end{enumerate}
Motifs 1--3 en phase~1 (comptage)~; 4--5 en automates finis (phase~2).
Onze commandes \texttt{user-*} sont parsées~; chaque entrée référence des
fichiers par chemin dépôt~\cite{perforce_p4d}. Tableau~\ref{tab:relisas}~:
mapping ReliSA (\emph{proxies}).

\subsection{Structure du journal \texttt{p4d}}
Le serveur Perforce journalise chaque requête cliente sous la forme
\texttt{Perforce server info:} suivie d'une ligne
horodatée~\cite{perforce_p4d}. Dans notre \texttt{log.txt} (1,7~Go),
chaque ligne utile commence par un horodatage, un \texttt{pid}, un
utilisateur \texttt{nom@workspace}, une adresse IP, un client entre
crochets (\texttt{[UE/v91]}, \texttt{[P4V/NTX64/...]}, etc.) et la
commande entre apostrophes. Les chemins de fichiers apparaissent dans
les arguments (\texttt{D:/...}, \texttt{C:/...}, parfois
\texttt{//depot/...}). La fig.~\ref{fig:log} montre un extrait
\textbf{simplifié et anonymisé} d'un \texttt{user-edit}~:
\textbf{(a)}~la commande~; \textbf{(b)}~la ligne \texttt{completed}
(même \texttt{pid}, durée en secondes)~; \textbf{(c)}~les métriques
serveur préfixées \texttt{---} (\texttt{lapse}, mémoire, RPC, accès
base). Le parseur assemble le bloc par \texttt{pid} uniquement~: dans
le journal réel, des milliers d'événements d'autres \texttt{pid}
s'intercalent entre \textbf{(a)} et \textbf{(b)} (ligne
\texttt{user-fstat} en exemple). Seules onze commandes \texttt{user-*}
sont retenues (edit, submit, sync, revert, \ldots)~; le reste est
ignoré.

\begin{figure}[!t]
  \centering
  \begin{minipage}{\linewidth}
\begin{lstlisting}[style=p4log]
Perforce server info:
	2026/04/26 22:34:38 pid 1000 auteur@workspace 10.0.0.1 [UE/v91] 'user-edit D:/depot/Projet/Content/Niveau/niveau.umap'
Perforce server info:
	2026/04/26 22:34:38 pid 2000 autre@client 10.0.0.2 [P4V/NTX64/...] 'user-fstat //workspace/...'
Perforce server info:
	2026/04/26 22:34:39 pid 1000 completed .032s
Perforce server info:
	2026/04/26 22:34:39 pid 1000 auteur@workspace 10.0.0.1 [UE/v91] 'user-edit D:/depot/Projet/Content/Niveau/niveau.umap'
--- lapse .032s
--- memory cmd/proc 10mb/10mb
--- rpc msgs/size in+out 0+1/0mb+0mb himarks 130372/64836 snd/rcv .000s/.000s
--- db.have
---   pages in+out+cached 1+0+2
\end{lstlisting}
  \end{minipage}
  \caption{Extrait simplifié et anonymisé du \texttt{log.txt} étudié
    (\texttt{user-edit}). \textbf{(a)}~commande~;
    \textbf{(b)}~\texttt{completed} (même \texttt{pid})~;
    \textbf{(c)}~métriques \texttt{---}. La ligne \texttt{user-fstat}
    illustre le bruit entremêlé, ignoré par le parseur.}
  \label{fig:log}
\end{figure}

\subsection{Architecture du pipeline}
Fig.~\ref{fig:arch} résume \toolname{}~: le journal brut
(\texttt{log.txt}, ${\sim}$1,7~Go) est lu \textbf{ligne par ligne} par
\texttt{LecteurLog}, sans chargement intégral en mémoire. Le parseur
retient onze commandes \texttt{user-*} et assemble, pour chacune, un
\textbf{bloc d'événement} en indexant les lignes par \texttt{pid}~:
commande $\rightarrow$ \texttt{completed} (durée) $\rightarrow$ métriques
\texttt{---} (fig.~\ref{fig:log}). Les blocs strictement identiques sont
ensuite retirés par déduplication
(\texttt{timestamp|pid|commande|args|fichiers})~; l'ordre chronologique
est conservé pour le traitement d'événements complexes. Le flux de blocs
dédupliqués alimente \beepbeep{}~3.13 en deux chaînes de processeurs~:
phase~1 (métriques OpenMetrics) et phase~2 (alertes CEP). Les sorties
sont écrites à la racine du projet~: dataset
\texttt{perforce\_events\_v3.json/.csv} (27 colonnes), quatre JSON de
métriques et \texttt{perforce\_events\_alertes.json}.

\figZeroArchitecture

\subsection{Phase~1~: m\'etriques d'observabilit\'e}
Counter, Gauge ($+1$/\texttt{edit}, $-1$/\texttt{submit}/\texttt{revert}),
Histogram (15 buckets), Summary (p50, p95, p99). Soutient procrastination,
Warm Bodies et Lone Wolf.

\subsection{Phase~2~: traitement d'\'ev\'enements complexes}
Automate de Moore par fichier~: \texttt{edit}$\rightarrow$ouvert,
\texttt{submit}$\rightarrow$fermé, \texttt{revert} depuis ouvert
$\rightarrow$alerte, ouvert en fin de flux $\rightarrow$apathie. Rush si
$\geq 3$ soumetteurs/heure. La jauge terminale doit égaler les alertes
d'apathie (section~\ref{sec:validation_croisee}). Alertes exportées en
JSON (fichier, utilisateur, horodatage, motif).

\section{R\'esultats}
\label{sec:resultats}

\textbf{Lecture des figures.}
Les figures illustrent l'infrastructure du pipeline
(fig.~\ref{fig:arch}, section~\ref{sec:methodo}), les métriques de
phase~1 (fig.~\ref{fig:compteurs} à~\ref{fig:quantiles_sync}) et trois
familles d'alertes de phase~2 (fig.~\ref{fig:alertes},
fig.~\ref{fig:rush}). Les antipatterns \emph{Warm Bodies} et
\emph{Lone Wolf} sont quantifiés en phase~1 et rapportés
numériquement (tableau~\ref{tab:lonewolf} pour Lone Wolf)~; une
visualisation dédiée pour ces deux motifs est laissée aux travaux
futurs.

Évaluation sur un journal de 1,7~Go issu d'un \textbf{cours}
multi-équipes (projet NAD/NAND, moteur de jeu 5.6, 26--27~avril~2026)---
contexte expérimental, distinct du déploiement industriel visé.
Le tableau~\ref{tab:data} résume l'extraction~; le flux complet est
traité en ${\sim}$14~s sur un ordinateur portable courant, ce qui
confirme que la conception en flux passe à l'échelle pour des journaux
de plusieurs gigaoctets. La fig.~\ref{fig:compteurs} donne une vue
d'ensemble des totaux phase~1 (événements, \texttt{edit},
\texttt{submit}).

\begin{table}[!t]
  \caption{Jeu de données et résumé de l'extraction.}
  \label{tab:data}
  \centering
  \scriptsize
  \begin{tabularx}{\linewidth}{@{}Xr@{}}
    \toprule
    \textbf{Quantité} & \textbf{Valeur} \\
    \midrule
    Taille du journal brut & 1,7~Go \\
    Événements extraits & 17\,187 \\
    Événements uniques (après dédup.) & 16\,262 \\
    Lignes dupliquées retirées & 925 \\
    Événements \texttt{edit} / \texttt{submit} & 4\,069 / 1\,224 \\
    Temps de traitement & ${\sim}$14~s \\
    \bottomrule
  \end{tabularx}
\end{table}

\figUnCompteurs

\subsection{Activité dans le temps}
Le taux de soumission présente la procrastination collective attendue
(fig.~\ref{fig:activite}, a.)~: quasi nul la nuit, puis un pic de
\textbf{107} soumissions à \textbf{15:00} (jour de remise). La jauge de
fichiers ouverts suit une dynamique corrélée (fig.~\ref{fig:activite},
b.), avec un pic à \textbf{500} fichiers ouverts et \textbf{79}
utilisateurs distincts en modification. Le rush collectif (phase~2)
atteint \textbf{60} soumetteurs distincts dans la même heure (27/04,
13~h, fig.~\ref{fig:rush}).

\figActivite

\figHuitRush

\subsection{Distribution des durées de commande}
La fig.~\ref{fig:durees} montre l'histogramme des durées (16\,232
observations, 15~buckets). Environ \textbf{89\,\%} des commandes
s'exécutent en moins de 0,25~s~; une minorité de \texttt{sync} sur
fichiers binaires forme une longue traîne (moyenne 7,36~s, médiane
0,094~s). Les fig.~\ref{fig:quantiles} et~\ref{fig:quantiles_sync}
détaillent les quantiles Summary par commande~; le
tableau~\ref{tab:quantiles} en donne les valeurs numériques.
\texttt{sync} atteint p99\,=\,1\,102~s.

\figQuatreHistogramme

\begin{table}[!t]
  \caption{Quantiles Summary (secondes).}
  \label{tab:quantiles}
  \centering
  \scriptsize
  \begin{tabularx}{\linewidth}{@{}Xrrr@{}}
    \toprule
    \textbf{Commande} & \textbf{p50} & \textbf{p95} & \textbf{p99} \\
    \midrule
    globale & 0,094 & 1,144 & 85,1 \\
    \texttt{edit} & 0,110 & 0,204 & 0,698 \\
    \texttt{submit} & 0,125 & 0,250 & 0,821 \\
    \texttt{sync} & 0,188 & 104,7 & 1\,102 \\
    \bottomrule
  \end{tabularx}
\end{table}

\figQuantilesCombine

\subsection{Antipatterns détectés}
La phase~2 signale \textbf{217} séquences
\emph{modification-puis-annulation} et \textbf{250} fichiers en
\emph{Bystander Apathy} (fig.~\ref{fig:alertes}). L'antipattern
\emph{Lone Wolf} fait ressortir plusieurs contributeurs extrêmes
(tableau~\ref{tab:lonewolf})~; le maximum atteint 364~modifications pour
2~soumissions ($\times$182). \emph{Warm Bodies} repose sur le seuil
$\theta_{\mathrm{wb}}$ en phase~1 (faible $E(u){+}S(u)$).

\figSeptAlertes

\begin{table}[!t]
  \caption{Ratios Lone Wolf les plus extrêmes (anonymisés).}
  \label{tab:lonewolf}
  \centering
  \scriptsize
  \begin{tabularx}{\linewidth}{@{}Xrrr@{}}
    \toprule
    \textbf{Utilisateur} & \textbf{Modif.} & \textbf{Soum.} & \textbf{Ratio} \\
    \midrule
    U1 & 364 & 2 & $\times$182 \\
    U2 & 230 & 2 & $\times$115 \\
    U3 & 216 & 3 & $\times$72 \\
    U4 & 90 & 2 & $\times$45 \\
    U5 & 140 & 4 & $\times$35 \\
    U6 & 85 & 3 & $\times$28 \\
    \bottomrule
  \end{tabularx}
\end{table}

\subsection{Validation croisée des deux phases}
\label{sec:validation_croisee}
Les deux phases sont volontairement redondantes sur une grandeur.
La jauge de phase~1 suit les fichiers ouverts via un solde $+1/-1$ et,
en fin de flux, rapporte \textbf{250} fichiers encore ouverts.
L'automate de phase~2, par un mécanisme indépendant (analyse fichier par
fichier), rapporte lui aussi exactement \textbf{250} alertes d'apathie.
L'accord est exact~: une erreur de comptage dans la jauge ou un défaut
dans l'automate ferait diverger ces deux valeurs. En l'absence de vérité
terrain, cette concordance constitue un indice interne fort de cohérence
entre les deux phases.

\section{Conclusion et limites}
\label{sec:conclu}

Nous avons présenté \toolname{}, un outil en flux qui détecte cinq
antipatterns de collaboration à partir des journaux serveur Perforce, en
combinant des métriques de style Prometheus (phase~1) et un détecteur
d'état sensible à l'ordre des événements (phase~2), orchestré par
\beepbeep{}. Sur un journal de cours de 1,7~Go, il traite
16\,262 événements en quelques secondes et produit à la fois un portrait
agrégé de l'activité et des alertes qu'un simple compteur sans mémoire
ne peut émettre. L'accord exact entre la jauge de fichiers ouverts et le
décompte d'apathie fondé sur l'automate (250/250,
section~\ref{sec:validation_croisee}) illustre la cohérence interne des
deux phases.

Les pics de soumission évoquent des remises de travail~; les
\texttt{sync} longues, un rattrapage d'assets binaires après absence~;
les ratios \texttt{edit}/\texttt{submit} élevés, un travail local peu
partagé. Ces interprétations restent toutefois \emph{hypothétiques} tant
qu'elles ne sont pas confirmées qualitativement auprès des équipes.

\textbf{Limites.}
L'évaluation repose sur \textbf{un seul journal} issu d'un cours
multi-équipes (26--27~avril~2026), contexte expérimental distinct d'un
déploiement industriel continu. Seuls \textbf{cinq comportements sur 168}
fiches ReliSA sont opérationnalisés, comme \emph{proxies} parfois
approximatifs (p.\,ex.\ edit$\rightarrow$revert pour
\emph{Cascading\_Branches}). Les \textbf{seuils} ($\rho$, $\tau$,
$\theta_{\mathrm{wb}}$, rush $\geq 3$) sont calibrés sur ce jeu de
données et ne sont pas généralisables sans réajustement. Le journal
mélange de \textbf{nombreux dépôts et équipes}~; les métriques globales
masquent des dynamiques locales. Les antipatterns \textbf{organisationnels}
(réunions, management, dette documentaire) sont \textbf{absents des
logs VCS}. \toolname{} ne voit que ce que le serveur enregistre~: un
\emph{Lone Wolf} peut se coordonner hors dépôt, et un fichier «~abandonné~»
peut correspondre à une branche exploratoire légitime. Les alertes sont
des \textbf{candidats à interprétation}, non des verdicts binaires.
La validation 250/250 (section~\ref{sec:validation_croisee}) ne couvre que
l'\textbf{état final du flux}~: un fichier encore ouvert à la fin du
journal n'équivaut pas à un abandon avéré. Le rapprochement
\texttt{edit}$\leftrightarrow$\texttt{submit} entre chemins locaux
(\texttt{D:/...}) et dépôt (\texttt{//...}) repose sur des
\textbf{heuristiques}. Les \texttt{user-submit -i}, ainsi que
\texttt{user-lock} et \texttt{user-shelve} (absents de notre log),
limitent le suivi fichier par fichier. L'analyse porte sur les
\textbf{métadonnées VCS} (commandes, chemins, durées), non sur le contenu
des assets binaires. Le pipeline est \textbf{hors ligne et par lots}~: il
ne reproduit pas un déploiement Prometheus temps réel. La
\textbf{déduplication} des blocs identiques peut masquer des répétitions
comportementales. L'outil est \textbf{spécifique à \texttt{p4d}}~; sa
transposition à Git ou à d'autres VCS reste à démontrer. Enfin, aucune
\textbf{mesure de précision} (vrais/faux positifs) n'a encore été
réalisée sur un échantillon d'alertes étiqueté.

\textbf{Perspectives.}
Valider empiriquement la précision des alertes (échantillon annoté,
estimation TP/FP), identifier des \emph{hotspots} de \texttt{revert} sur
fichiers binaires, produire une synthèse anonymisée par équipe, et
publier le code d'\toolname{} sur un dépôt GitHub pour favoriser la
reproductibilité.

\bibliographystyle{IEEEtran}
\bibliography{ieee_gem_references}

\end{document}
